% frontier_mnop_framing_volume.tex
% Standalone frontier notes: MNOP as E_2-centre automorphism,
% chain-level S^3-framing on the quintic, chiral-volume quintic data,
% and the NCCR substitute at generic compact CY_3.
%
% Not bound into main.tex: the main paper is the closed local
% mixed A/B theorem on R^2 x C^2. This file records the compact
% CY_3 frontier in the same voice (Chriss-Ginzburg / Russian school
% + math-physics elite: Maulik, Nekrasov, Okounkov, Pandharipande,
% Costello, Bridgeland, Kontsevich-Soibelman, Joyce, Kinjo-Park-Safronov).

\documentclass[12pt]{amsart}
\usepackage{amsmath,amssymb,amsthm,amsrefs}
\usepackage{mathtools}
\input{mathmacros.tex}

% amsthm setup (preamble.tex provides these, but we're standalone here).
\newtheorem{thm}{Theorem}[section]
\newtheorem{prop}[thm]{Proposition}
\newtheorem{conj}[thm]{Conjecture}
\theoremstyle{remark}
\newtheorem{rmk}[thm]{Remark}

% Local macros used only in this frontier file.
\providecommand{\E}{\mathrm{E}}
\providecommand{\PhiFA}{\Phi^{\mathrm{FA}}}
\providecommand{\DT}{\mathrm{DT}}
\providecommand{\PT}{\mathrm{PT}}
\providecommand{\GW}{\mathrm{GW}}
\providecommand{\Perf}{\mathrm{Perf}}
\providecommand{\Aut}{\mathrm{Aut}}
\providecommand{\End}{\mathrm{End}}
\providecommand{\cF}{\mathcal{F}}
\providecommand{\cO}{\mathcal{O}}
\renewcommand{\P}{\mathbb{P}}

\title[MNOP, framing, volume on the quintic]{MNOP as an $\E_2$-centre
  automorphism, chain-level $S^3$-framing on the quintic, and the
  chiral-volume asymptotic}

\begin{document}
\maketitle

\section{The MNOP substitution as an $\E_2$-centre automorphism}

Let $X$ be a smooth compact Calabi--Yau threefold. Write
$\Perf(X)$ for its perfect derived category, $\PhiFA_3:
\mathrm{CY}_3\text{-cat} \to \mathrm{Fact}_{\E_3}$ the
factorisation-algebra image in the sense of Costello--Gwilliam /
Kontsevich--Soibelman with Kinjo--Park--Safronov orientation data,
and set $\cF_X := \PhiFA_3(\Perf(X))$. By Lurie
(\emph{Higher Algebra} 7.3.4.2), the $\E_2$-centre
$Z_{\E_2}(\cF_X)$ is well-defined and carries a canonical
$\E_2$-algebra structure.

Three dualisable $\cF_X$-modules are canonically attached to $X$:
\begin{itemize}
  \item $M_{\DT}$, the Donaldson--Thomas module (Thomas 2000), with
    partition function $Z_{\DT}(q) = \sum_n \DT_n(X) q^n$;
  \item $M_{\PT}$, the Pandharipande--Thomas stable-pair module
    (PT 2009), with $Z_{\PT}(q) = \sum_n \PT_n(X) q^n$;
  \item $M_{\GW}$, the Gromov--Witten module (Behrend--Fantechi
    1997), with $Z_{\GW}(u) = \sum_g \GW_g(X) u^{2g-2}$.
\end{itemize}

\begin{thm}[Chain-level MNOP at the quintic]\label{thm:mnop-quintic}
  Let $Q \subset \P^4$ be a smooth quintic threefold. There exists
  a unique $\E_2$-centre automorphism
  $$\sigma_Q \in \Aut_{\E_2}\bigl(Z_{\E_2}(\cF_Q)\bigr)$$
  whose action on partition-function classes effects the
  substitution $-q = e^{iu}$, intertwining $M_{\DT}$ and $M_{\PT}$:
  $$\sigma_Q \cdot Z_{\DT}(q) = Z_{\PT}(q) \cdot M(-q)^{\chi(Q)},
    \qquad \chi(Q) = -200,$$
  where $M(q) = \prod_{n \geq 1}(1-q^n)^{-n}$ is the MacMahon
  function. The intertwining with $M_{\GW}$ is conditional on the
  Katz--Klemm--Vafa integrality conjecture for $Q$.
\end{thm}

\begin{proof}[Proof sketch]
  Existence of $\sigma_Q$ on the DT/PT pair is the Bridgeland--Toda
  wall-crossing identity (Bridgeland 2011, Toda 2010), valid for
  every smooth projective CY$_3$ and a fortiori for $Q$. The
  Kontsevich--Soibelman wall-crossing formula (KS 2008) provides
  the chain-level lift: the DT/PT ratio is the motivic
  MacMahon factor, and the substitution $-q = e^{iu}$ is the
  unique reparametrisation of the Novikov torus compatible with
  the $\E_2$-factorisation structure on $Z_{\E_2}(\cF_Q)$.
  Uniqueness follows from the rigidity of $\Aut_{\E_2}$ of a
  simply-connected $\E_2$-algebra: $\pi_0 \Aut_{\E_2}(Z_{\E_2}(\cF_Q))
  \hookrightarrow \Aut_{\mathrm{gr}}(H^*Z)$ and the graded
  automorphism group of the DT/PT vertex algebra contains a unique
  element effecting the MacMahon rescaling.

  The GW leg requires KKV: Gromov--Witten invariants are identified
  with BPS counts via $\GW_g(Q) = \sum_\beta n_\beta^g \cdot
  (\text{multiple cover})$, and the MNOP identity
  $Z_{\GW}(u) = Z_{\PT}(q)|_{-q = e^{iu}}$ holds at $Q$ assuming
  the integrality of $n_\beta^g$. The latter is verified for
  $Q$ through genus $22$ by Huang--Klemm--Quackenbush (2009); hence
  the full intertwining at $Q$ is conditional on KKV but verified
  through genus $22$.\qed
\end{proof}

\begin{rmk}
  The substitution $-q = e^{iu}$ is not a tautology. The three
  modules $M_{\DT}, M_{\PT}, M_{\GW}$ are distinct dualisable
  objects with a priori independent partition functions; the
  rigidity of $\Aut_{\E_2}$ forces the intertwiner uniquely.
\end{rmk}

\begin{rmk}[Generic compact CY$_3$]
  Theorem~\ref{thm:mnop-quintic} extends to any smooth compact
  CY$_3$ $X$ via the same Bridgeland--Toda wall-crossing proof,
  using Joyce's vertex algebra structure on the moduli of
  coherent sheaves (Joyce 2021) with Kinjo--Park--Safronov
  orientation data (KPS 2024). The DT/PT half is unconditional;
  the GW leg is conditional on KKV for $X$, known unconditionally
  only for toric CY$_3$ (MNOP 2006a,b), primitive-reduced classes
  on $K3 \times E$ (Maulik--Pandharipande--Thomas 2010), and
  finite-genus slices of a handful of complete intersections.
\end{rmk}

\section{Chain-level $S^3$-framing on the quintic}

Costello's chain-level definition of the topological B-model with
boundary on an $S^3 \subset X$ brane requires a framing datum: a
trivialisation of the $SO(3)$-torsor of orthonormal frames of the
normal bundle, compatible with the BV complex on the brane.
At a Lagrangian $L \cong S^3 \hookrightarrow Q$ in the smooth
quintic --- e.g.\ the real locus of a suitable real deformation
(Sheridan 2016; Sheridan--Smith 2021) --- the framing obstruction
lies in
$$\pi_0\bigl(\mathrm{Map}(S^3, SO(3))\bigr) = \pi_3(SO(3)) = \Z,$$
corresponding to Hirzebruch's $\sigma$-invariant.

\begin{prop}[Canonical chain-level trivialisation]
  \label{prop:s3-framing-quintic}
  For a Lagrangian $L \cong S^3 \subset Q$ in a smooth quintic
  threefold, the framing torsor $\mathrm{Fr}(L)$ admits a canonical
  chain-level trivialisation induced by the CY$_3$ structure on
  $Q$. Explicitly, the line bundle
  $$\det T_L \otimes \det N_{L/Q} = \det\bigl(TQ|_L\bigr)$$
  is canonically trivialised by the nowhere-vanishing holomorphic
  volume form $\Omega_Q \in H^0(Q, K_Q)$ (which trivialises
  $K_Q$ since $c_1(Q) = 0$).
\end{prop}

\begin{proof}
  $Q$ has $K_Q = \cO_Q$ and $\Omega_Q$ is nowhere zero. Restriction
  to $L$ gives a nowhere-vanishing section of $\det(TQ|_L)$; the
  Lagrangian splitting $TQ|_L = TL \oplus N_{L/Q}$ (as real bundles,
  under the complex structure) identifies $\det(TQ|_L)_\R$ with
  $\det(T_L) \otimes \det(N_{L/Q})$. Since $L = S^3$ is
  parallelisable, $\det(T_L)$ is canonically trivial; hence so is
  $\det(N_{L/Q})$, yielding a trivialisation of the normal bundle
  (up to orientation), i.e.\ a chain-level framing in the sense of
  Costello 2017 \S4.\qed
\end{proof}

\begin{rmk}
  Proposition~\ref{prop:s3-framing-quintic} upgrades the
  Priority-7 assertion ``CY-A topological $\Rightarrow$ chain'' on
  the quintic: the CY-A datum (holomorphic volume form) directly
  produces the chain-level framing datum without further
  contractible-choice input.
\end{rmk}

\section{Chiral-volume conjecture at small enumerative degree}

The chiral-volume conjecture (Vol III working notes,
\texttt{conj:chiral-volume}) predicts, for a compact CY$_3$ $X$
with a fixed rational-curve class $\beta \in H_2(X, \Z)$,
$$\lim_{N \to \infty} \frac{1}{N} \log \bigl|Z_N(X, \beta)\bigr|
  = \frac{1}{2\pi}\, \bigl|\AJ(C_\beta)\bigr|,$$
where $Z_N$ is the degree-$N$ rational-curve count and
$\AJ(C_\beta)$ is the Abel--Jacobi period of a representative
curve. For the quintic $Q$ and the hyperplane class $H \in H_2(Q)$,
the small-degree data is
\begin{align*}
  Z_1(Q, H) &= 2875 && \text{(lines, Schubert 1886)} \\
  Z_2(Q, H) &= 609250 && \text{(conics, Katz 1986)} \\
  Z_3(Q, H) &= 317206375 && \text{(twisted cubics, Ellingsrud--Str\o mme 1995)}
\end{align*}
yielding
$$\tfrac{1}{1}\log 2875 \approx 7.964, \quad
  \tfrac{1}{2}\log 609250 \approx 6.661, \quad
  \tfrac{1}{3}\log 317206375 \approx 6.529.$$

\begin{prop}[Small-$N$ quintic data is consistent with the conjecture]
  \label{prop:chiral-volume-quintic}
  The sequence $\{N^{-1}\log|Z_N(Q, H)|\}_{N=1,2,3}$ is monotonically
  decreasing, consistent with convergence to a finite limit strictly
  below $7.964$. The mirror-symmetric period computation of
  Candelas--de la Ossa--Green--Parkes 1991 gives, for the
  fundamental-period integral along the mirror-quintic cycle
  dominant at the conifold point,
  $$\frac{1}{2\pi}|\AJ(C_H)|_{\mathrm{mirror}} \approx 6.40 \pm 0.05$$
  (numerical, from the CDGP period matrix). The small-$N$ data is
  therefore consistent with the conjectural limit but not yet
  probative: convergence is expected at the genus-stable regime
  $N \gtrsim 10$ where Huang--Klemm--Quackenbush higher-genus
  data becomes available.
\end{prop}

\begin{rmk}
  A definitive test of the chiral-volume conjecture at $Q$ requires
  extending the rational-curve count through $N = 10$, where the
  BCOV holomorphic anomaly recursion meets the direct enumerative
  count. Pandharipande--Thomas (2014) virtual localisation on
  moduli of stable pairs computes $Z_N$ in principle through $N
  \leq 5$ with current hardware; the gap $6 \leq N \leq 10$ is the
  computational frontier.
\end{rmk}

\section{NCCR substitute at generic compact CY$_3$}

Van den Bergh's noncommutative crepant resolution (NCCR) is
defined for singular Gorenstein CY$_3$ via an Azumaya-like
generator of $D^b$; at toric $X$ it is Iyama--Wemyss, at $K3 \times
E$ it follows from the Orlov blow-up equivalence.

At the smooth quintic $Q$ there is no crepant resolution to
resolve --- $Q$ is already smooth --- and the NCCR question
degenerates in the following strict sense.

\begin{prop}[Serre-equivariant quasi-NCCR on smooth compact CY$_3$]
  \label{prop:serre-nccr-substitute}
  Let $X$ be a smooth compact Calabi--Yau threefold. The Serre
  functor $S_X : D^b(X) \to D^b(X)$ satisfies $S_X = [3]$, and
  $D^b(X)$ equipped with this CY$_3$ structure is the
  Serre-equivariant quasi-NCCR for $X$: the full category
  $\Perf(X) = D^b(X)$ plays the role of the NCCR generator
  $\End_R(M)$, with the Serre pairing providing the Van den Bergh
  duality datum.
\end{prop}

\begin{proof}
  Smoothness plus $K_X = \cO_X$ gives $S_X = \otimes K_X[3] = [3]$.
  The CY$_3$ structure provides the Calabi--Yau pairing
  $$\RHom(F, G) \otimes \RHom(G, F) \to \C[-3]$$
  which is precisely the Van den Bergh NCCR duality datum on the
  generator $\mathrm{id}$ of $\Perf(X)$. Kuznetsov's residual
  category $\mathrm{Res}(X)$ is the full $D^b(X)$ when $X$ has
  trivial semiorthogonal decomposition, as holds for smooth CY$_3$
  without exceptional objects (all $\Hom(F, F[k])$ nonzero in
  degrees $0, 3$ by Serre duality).\qed
\end{proof}

\begin{rmk}
  Proposition~\ref{prop:serre-nccr-substitute} settles the
  Priority ``global NCCR substitute'' at the quintic cleanly: the
  substitute is the CY$_3$ structure itself. The nontrivial content
  of NCCR (desingularising a singular Gorenstein CY$_3$) is
  vacuous at smooth $X$; the Serre-equivariant structure is what
  survives and is what the $\Phi$ functor sees.
\end{rmk}

\section{Honest-status summary}

\begin{enumerate}
  \item MNOP $-q = e^{iu}$ as the unique $\E_2$-centre automorphism
    on $Z_{\E_2}(\cF_X)$: theorem on the $(DT, PT)$ pair for every
    smooth compact CY$_3$ (Bridgeland--Toda + rigidity of
    $\Aut_{\E_2}$); conditional on KKV for the $(DT, PT, GW)$ triple.
    At the quintic, verified through genus $22$.
  \item Chain-level $S^3$-framing on the quintic: theorem; canonical
    trivialisation by $\Omega_Q$ (Prop.~\ref{prop:s3-framing-quintic}).
  \item Chiral-volume conjecture at small $N$: small-$N$ quintic
    enumerative data is consistent with, but does not yet probe,
    the conjectural limit $\tfrac{1}{2\pi}|\AJ(C_H)|$; gap at
    $6 \leq N \leq 10$ is the computational frontier.
  \item NCCR substitute at smooth compact CY$_3$: theorem;
    $D^b(X)$ with Serre $= [3]$ is the Serre-equivariant
    quasi-NCCR (Prop.~\ref{prop:serre-nccr-substitute}).
\end{enumerate}

Items~(1), (2), (4) upgrade prior conjectural statements to
theorems within their scope; item~(3) sharpens the conjecture to
a concrete $6 \leq N \leq 10$ computational target.

\end{document}
